\section{Framens fält}
\label{sec:frame}

Kursens protokoll använder följande struktur. Siffran är fältets längd i bytes:

\begin{console}
SOF (2) | LEN (1) | TYPE (1) | DST (1) | SRC (1) | SEQ (2) | DATA (N) | CHK (2)
\end{console}

\begin{booktable}{@{}llL@{}}
\thead{Fält} & \thead{Bytes} & \thead{Vad det löser} \\ \midrule
SOF & 2 & \emph{Var börjar meddelandet?} Ett fast mönster, \code{0xA5F7}, som markerar start. \\
LEN & 1 & \emph{Hur långt är det?} Antalet databytes som följer. \\
TYPE & 1 & \emph{Vad betyder det?} Meddelandets sort. \\
DST & 1 & \emph{Till vem?} Mottagarens adress. \\
SRC & 1 & \emph{Från vem?} Avsändarens adress. \\
SEQ & 2 & Löpnummer, som gör det möjligt att para ihop svar med fråga och att känna igen en
  dubblett. \\
DATA & N & Nyttolasten, \code{LEN} bytes lång. \\
CHK & 2 & \emph{Kom det fram helt?} Summan av alla föregående bytes. \\
\end{booktable}

Frametyperna i kursens protokoll:

\begin{narrowtable}{@{}ll@{}}
\thead{Typ} & \thead{Värde} \\ \midrule
\code{Ping} & \code{0x00} \\
\code{Pong} & \code{0x01} \\
\code{StatusRequest} & \code{0x02} \\
\code{StatusResponse} & \code{0x03} \\
\code{Ack} & \code{0x04} \\
\code{Nack} & \code{0x05} \\
\end{narrowtable}

\subsection{Varför SOF är två bytes}

En byte räcker inte. Vilket värde man än väljer kommer det förr eller senare att dyka upp mitt i
en nyttolast, och då hittar mottagaren en frame som inte finns.

Två bytes gör det osannolikt snarare än omöjligt --- \code{0xA5F7} kan fortfarande förekomma i
data. Det är därför SOF ensamt aldrig är tillräckligt: mottagaren måste också kontrollera att
längden är rimlig och att checksumman stämmer innan den tror på framen. \Kapref{ch:parsning}
bygger den kontrollen.

\subsection{Varför LEN behövs när DATA ändå tar slut}

Därför att mottagaren läser en byte i taget och inte kan se framåt. Utan ett längdfält vet den
aldrig om nästa byte är data eller checksumma.

Alternativet vore ett slutmärke, men då måste det mönstret vara omöjligt i nyttolasten --- vilket
kräver att man kodar om datan. CAN löser samma problem med bitstoppning (se kursens CAN-bok);
här räcker ett längdfält.

\subsection{Checksumman}

Kursens checksumma är enkel: \textbf{16-bitars summering av alla bytes fram till CHK-fältet}.

\begin{cppcode}
std::uint16_t computeChecksum(const std::uint8_t* buf, const std::size_t chkOffset) noexcept
{
    std::uint16_t checksum{};
    for (std::size_t i{}; i < chkOffset; ++i)
    {
        checksum += buf[i];
    }
    return checksum;
}
\end{cppcode}

Den upptäcker enstaka bitfel och de flesta korta felburstar. Den upptäcker däremot inte att två
bytes bytt plats, och inte att en byte ökat med ett medan en annan minskat med ett --- summan är
densamma. En CRC gör det bättre, och det är vad CAN använder.

\begin{aside}
\note Checksumman säger att framen kom fram \emph{hel}. Den säger ingenting om att den kom fram
\emph{alls}. Skillnaden mellan integritet och leveransgaranti är hela \kapref{ch:bussar}s andra
hälft, och den är värd att hålla isär från början.
\end{aside}

\subsection{Serialisering och endianness}

Att serialisera är att skriva ut framens fält i rätt ordning i en buffer. Fält som är bredare än
en byte --- SOF, SEQ och CHK --- måste skrivas i en bestämd byteordning, och kursens protokoll
använder \textbf{big-endian}: den mest signifikanta byten först.

\begin{console}
SEQ = 0x7F05  ->  7F 05      (big-endian, det protokollet säger)
              ->  05 7F      (little-endian, det din x86 gör i minnet)
\end{console}

Det här är den klassiska buggen i all protokollkod, och den är obehaglig av två skäl. Den är
osynlig när båda sidor har samma fel, och den ser ut som slumpmässig korruption när de inte har
det: \code{0x0005} blir \code{0x0500}, alltså 1280 i stället för 5.

Skriv därför aldrig ut en flerbytesvariabel genom att kopiera minnet. Skriv ut den byte för byte,
med explicita skiftningar:

\begin{cppcode}
// Big-endian: most significant byte first, whatever the host's byte order is.
buf[offset + 0U] = static_cast<std::uint8_t>(value >> 8U);
buf[offset + 1U] = static_cast<std::uint8_t>(value);
\end{cppcode}

\subsection{Ett komplett exempel}

En \code{StatusRequest} från adress \code{0x17} till adress \code{0x25}, med sekvensnummer
\code{0x7F05} och utan nyttolast:

\begin{console}
A5 F7 00 02 25 17 7F 05 02 5E
|---| |  |  |  |  |---| |---|
SOF  LEN TYPE      SEQ   CHK
        DST--^  ^--SRC
\end{console}

Checksumman: $\mathtt{A5} + \mathtt{F7} + \mathtt{00} + \mathtt{02} + \mathtt{25} + \mathtt{17} +
\mathtt{7F} + \mathtt{05} = 606 = \mathtt{0x025E}$.

Notera att \code{DST} står före \code{SRC}. Ordningen är godtycklig men måste vara densamma på
båda sidor, och att blanda ihop dem ger ett system där varje nod envist svarar sig själv.
